\chapter{Evaluation, Testing, and Discussion}

This chapter assesses the effectiveness of the implementation described in Chapter~5.
It covers automated test coverage, manual campaign evaluation on live streaming targets,
prompt evolution metrics across two concrete use cases, tool description engineering
results, and a structured discussion of the platform's strengths and remaining
limitations.

\section{Evaluation approach}

Evaluation was performed with two complementary tracks:

\begin{itemize}
  \item automated checks (routing, tool behavior, persistence, regression safety);
  \item manual campaign evaluation on live matches and real streaming websites.
\end{itemize}

Manual evaluation was necessary because difficult failure modes (popups, nested iframes,
challenge screens, delayed manifests) are highly context-dependent and not fully
represented by static fixtures.

\section{Manual live-match evaluation protocol}

The manual protocol used during evaluation is illustrated in
Figure~\ref{fig:eval-protocol}.

\begin{figure}[H]
\centering
\begin{tikzpicture}[node distance=0.5cm]

\node[flowbox, minimum width=9cm, fill=blue!10, draw=blue!60] (sel)
  {Select live or actively streaming target URLs};
\node[flowbox, minimum width=9cm, below=0.45cm of sel] (run)
  {Run full pipeline (classification $\to$ specialist $\to$ harvest)};
\node[flowbox, minimum width=9cm, below=0.45cm of run] (class)
  {Record classification path and routing decision};
\node[flowbox, minimum width=9cm, below=0.45cm of class] (trace)
  {Inspect screenshots and tool traces after each major step};
\node[flowbox, minimum width=9cm, below=0.45cm of trace] (verify)
  {Verify extracted stream evidence and fallback behavior};

\node[decisionbox, minimum width=5cm, below=0.65cm of verify] (ok)
  {Evidence\\valid?};

% Success box: placed right of diamond but shifted down so it stays in margin
\node[flowbox, minimum width=4cm, fill=green!10, draw=green!60,
      right=1.8cm of ok] (pass)
  {Record\\as success};

% Failure box: placed below diamond, left-aligned
\node[flowbox, minimum width=4cm, fill=red!10, draw=red!60,
      left=1.8cm of ok] (fail)
  {Classify failure mode};

\node[flowbox, minimum width=9cm, fill=orange!8,
      below=2.0cm of ok] (improve)
  {Feed findings into prompt or tool improvement cycle};

% Labels under fail box
\node[font=\tiny, text=gray!70, below=0.1cm of fail, align=center]
  {popup $|$ iframe $|$ token $|$ anti-bot};

\foreach \a/\b in {sel/run, run/class, class/trace, trace/verify, verify/ok}
  \draw[myarrow] (\a) -- (\b);
\draw[myarrow] (ok.east) -- node[above,font=\tiny]{Yes} (pass.west);
\draw[myarrow] (ok.west) -- node[above,font=\tiny]{No}  (fail.east);
\draw[myarrow] (pass.south) -- ++(0,-0.5) -| (improve.east);
\draw[myarrow] (fail.south) -- ++(0,-0.9) -| (improve.west);

\end{tikzpicture}
\caption{Manual live-match evaluation protocol flowchart.}
\label{fig:eval-protocol}
\end{figure}

\section{Evaluation dimensions}

\begin{table}[H]
\centering
\begin{tabularx}{\textwidth}{>{\raggedright\arraybackslash\bfseries}p{3cm}
                              >{\raggedright\arraybackslash}X
                              >{\raggedright\arraybackslash}X}
\toprule
\rowcolor{accent!8}
Dimension & What is measured & Why it matters \\
\midrule
Routing quality & correctness of page-type and specialist selection &
  misrouting invalidates downstream extraction \\
Extraction quality & stream evidence completeness and validity &
  core objective of the platform \\
Fallback quality & hosting to embedded recovery success &
  measures platform recovery depth \\
Runtime reliability & stability under UI drift and anti-bot pressure &
  indicates operational robustness \\
Cost and cache behavior & token usage, cache ratio, model impact &
  determines scalability and sustainability \\
\bottomrule
\end{tabularx}
\caption{Evaluation dimensions used for this project.}
\label{tab:eval-dimensions}
\end{table}

\section{Prompt evaluation: navigation use case}

\subsection{Scenario description}

The test site presents a listing page with three tab filters: \emph{Live Now},
\emph{Today}, and \emph{Tomorrow}. Only the active tab renders match cards in the DOM;
the others are hidden. Clicking a tab triggers a JavaScript fetch that inserts new
cards without a full page reload.

\subsection{Evaluation methodology}

Each prompt version was run 10 times on the same target site during a live event window
(at least 5 matches visible across tabs). The metric was \emph{candidate URL recall}:
the ratio of discovered hosting URLs to total actual hosting URLs across all tabs.

\subsection{Results}

\begin{table}[H]
\centering
\begin{tabularx}{\textwidth}{>{\raggedright\arraybackslash\bfseries}p{3.6cm}
                              >{\raggedright\arraybackslash}X
                              >{\raggedright\arraybackslash}X
                              >{\raggedright\arraybackslash}X}
\toprule
\rowcolor{accent!8}
Prompt version & Avg. candidates found & Actual total & Recall \\
\midrule
Intent-based (v1) & 3.2 & 8 & 40\% \\
ReAct (v2) & 6.8 & 8 & 85\% \\
ReAct + budget cap (v3) & 7.1 & 8 & 89\% \\
\bottomrule
\end{tabularx}
\caption{Navigation use case: candidate URL recall by prompt version.}
\label{tab:nav-recall}
\end{table}

\textbf{Key finding.} The intent-based prompt consistently missed tab-gated content
because it had no instruction to look for navigational controls. The ReAct prompt's
explicit THINK step for detecting tabs more than doubled recall. The budget cap (v3)
added minimal extra coverage while preventing runaway tool calls on sites with many
empty tabs.

\subsection{Failure pattern analysis}

The remaining 11\% gap in v3 recall came from two patterns:
\begin{itemize}
  \item matches added dynamically within the observation window after the agent had
    already finished checking a tab;
  \item tabs that required a secondary scroll to reveal cards below the fold (caught by
    a separate scroll-and-reinspect loop introduced in v4).
\end{itemize}

\section{Prompt evaluation: player activation use case}

\subsection{Scenario A: AlbaPlayer on host page}

\textbf{Setup.} An AlbaPlayer instance embedded in an iframe within a host page. The
player shows a play triangle but does not auto-start. Clicking play triggers an HLS
manifest request.

\begin{table}[H]
\centering
\begin{tabularx}{\textwidth}{>{\raggedright\arraybackslash\bfseries}p{3.6cm}
                              >{\raggedright\arraybackslash}X
                              >{\raggedright\arraybackslash}X}
\toprule
\rowcolor{accent!8}
Prompt version & Behavior observed & m3u8 recovery rate \\
\midrule
Intent-based (v1) & Harvested without clicking play; 0 streams returned &
  0\% (10/10 failures) \\
ReAct (v2) & Called get\_media\_state, detected idle, clicked play, then harvested &
  70\% (7/10 successes) \\
ReAct + popup handler (v3) & v2 plus auto-close of new tab; js\_click fallback &
  90\% (9/10 successes) \\
\bottomrule
\end{tabularx}
\caption{Player activation (AlbaPlayer): m3u8 recovery rate by prompt version.}
\label{tab:alba-recovery}
\end{table}

\textbf{Key finding.} The v1 failure was 100\% systematic: no prompt instruction
directed the agent to check player state before harvesting. Adding \texttt{get\_media\_state}
as a mandatory pre-harvest step in v2 immediately recovered the majority of cases.
The remaining v2 failures were all caused by ad popups opening on the click event ---
resolved by the popup handler in v3.

\subsection{Scenario B: Tokenized Clappr player on embedded page}

\textbf{Setup.} A Clappr player on an embedded page where clicking play triggers a
JavaScript fetch that exchanges a token and then initiates the HLS manifest request.
The manifest URL contains a short-lived signed token as a query parameter.

\begin{table}[H]
\centering
\begin{tabularx}{\textwidth}{>{\raggedright\arraybackslash\bfseries}p{3.6cm}
                              >{\raggedright\arraybackslash}X
                              >{\raggedright\arraybackslash}X}
\toprule
\rowcolor{accent!8}
Prompt version & Behavior observed & m3u8 recovery rate \\
\midrule
ReAct v2 & Clicked play, harvested immediately with 1s wait; token not yet issued &
  40\% \\
ReAct v3 & Clicked play, harvested with 2s wait; caught majority of token exchanges &
  75\% \\
ReAct v3 + perf API & Also scanned performance.getEntriesByType for retroactive
  manifest entries & 85\% \\
\bottomrule
\end{tabularx}
\caption{Tokenized player (Clappr): m3u8 recovery rate under signed token flow.}
\label{tab:token-recovery}
\end{table}

\textbf{Key finding.} Token generation latency varies by site. A fixed 2s wait was
better than 1s but still missed fast-rotating tokens. Adding a retroactive scan via
the browser Performance API caught manifests that had already been requested before
the agent called harvest.

\section{Tool use evaluation}

\subsection{What was measured}

Tool use evaluation inspected the agent's tool call logs to answer three questions:
\begin{enumerate}
  \item Did the agent call the \emph{right} tool for each step (tool selection accuracy)?
  \item Did the agent call tools in the \emph{right order} (tool sequencing)?
  \item Did the agent call \emph{too many} tools before giving up (tool budget compliance)?
\end{enumerate}

\subsection{Common tool misuse patterns observed}

\begin{table}[H]
\centering
\begin{tabularx}{\textwidth}{>{\raggedright\arraybackslash\bfseries}p{3.8cm}
                              >{\raggedright\arraybackslash}X
                              >{\raggedright\arraybackslash}X}
\toprule
\rowcolor{accent!8}
Misuse pattern & Root cause & Fix applied \\
\midrule
\texttt{harvest()} before play click & Agent description for harvest did not warn about
  premature use & Added ``Use only after player activation confirmed'' to description \\
\texttt{inspect\_page} instead of \texttt{inspect\_hosting} & Agent had both tools;
  generic description did not signal specialization & Added ``Use instead of
  inspect\_page when on a host page'' \\
\texttt{screenshot()} called 5+ times in sequence & Agent used screenshot as a
  wait mechanism & Added explicit wait parameter to \texttt{interact} and removed
  screenshot-as-wait pattern from prompt \\
\texttt{navigate()} on a same-page tab click & Agent confused tab clicks
  (JS event) with navigation (full load) & Clarified in interact description:
  tab clicks are interact, not navigate \\
\bottomrule
\end{tabularx}
\caption{Common tool misuse patterns discovered through tool call log analysis.}
\label{tab:tool-misuse}
\end{table}

\subsection{Tool description improvements and measured effect}

After each tool description update, a before/after comparison was run on the same
set of 20 URLs. Table~\ref{tab:desc-effect} summarizes the impact.

\begin{table}[H]
\centering
\begin{tabularx}{\textwidth}{>{\raggedright\arraybackslash\bfseries}p{4cm}
                              >{\centering\arraybackslash}p{2.8cm}
                              >{\centering\arraybackslash}X}
\toprule
\rowcolor{accent!8}
Description change & Misuse rate before & Misuse rate after \\
\midrule
Added activation guard to \texttt{harvest} & 65\% premature harvest & 12\% \\
Added specialization signal to \texttt{inspect\_hosting} & 40\% wrong inspect & 8\% \\
Clarified \texttt{navigate} vs \texttt{interact} & 30\% nav/interact confusion & 5\% \\
\bottomrule
\end{tabularx}
\caption{Effect of tool description improvements on observed misuse rates.}
\label{tab:desc-effect}
\end{table}

\subsection{Token-based evaluation}

Beyond success rate, every run was analyzed for token efficiency. The key metrics
tracked were:

\begin{itemize}
  \item \textbf{Total tokens per run:} sum of input + output tokens across all LLM
    calls in the run;
  \item \textbf{Cache hit ratio:} fraction of input tokens served from the provider
    cache rather than recomputed;
  \item \textbf{Tokens per successful extraction:} total run tokens divided by number
    of confirmed m3u8 URLs found;
  \item \textbf{Tool call count:} total number of tool invocations per run (proxy for
    agent loop depth).
\end{itemize}

\begin{table}[H]
\centering
\begin{tabularx}{\textwidth}{>{\raggedright\arraybackslash\bfseries}p{3.8cm}
                              >{\raggedright\arraybackslash}X
                              >{\raggedright\arraybackslash}X}
\toprule
\rowcolor{accent!8}
Phase & Avg. tokens & Avg. tool calls \\
\midrule
Classification & \textasciitilde{}12k & 2--3 \\
Landing exploration & \textasciitilde{}45k & 4--8 \\
Hosting extraction & \textasciitilde{}180k & 8--14 \\
Embedded fallback & \textasciitilde{}110k & 6--10 \\
Infrastructure enrichment & \textasciitilde{}8k & 1--2 \\
\bottomrule
\end{tabularx}
\caption{Average token consumption and tool call count by run phase.}
\label{tab:token-by-phase}
\end{table}

\textbf{Key finding.} Hosting extraction dominates cost because it loops across
multiple server attempts, each requiring inspect $\to$ interact $\to$ screenshot $\to$
harvest. The most effective cost reduction was raising the confidence threshold for
early success exit: if \texttt{harvest} returns a confirmed m3u8 on the first server,
the agent exits immediately without trying remaining servers.

\section{Prompt iteration evaluation}

Each significant prompt change was evaluated against a fixed reference set of 20 URLs
covering both site families. The evaluation recorded: stream recovery rate, navigation
recall, tool call count, and first-server success rate.

\begin{table}[H]
\centering
\begin{tabularx}{\textwidth}{>{\raggedright\arraybackslash\bfseries}p{1.5cm}
                              >{\raggedright\arraybackslash}X
                              >{\raggedright\arraybackslash}X
                              >{\centering\arraybackslash}p{2.2cm}}
\toprule
\rowcolor{accent!8}
Version & Key change & Observable effect & Avg. tool calls \\
\midrule
v1 & Intent-based; no loop instruction & Early stops, no tab navigation & 3.1 \\
v2 & ReAct loop instruction; tab-check added & Tab-gated matches discovered & 6.2 \\
v3 & Popup handler; js\_click fallback added & Popup-blocked pages recovered & 7.8 \\
v4 & \texttt{get\_media\_state} before harvest & Premature harvest eliminated & 8.4 \\
v5 & Tool budget cap; early success exit & Cost reduced; recovery maintained & 7.1 \\
\bottomrule
\end{tabularx}
\caption{Prompt version history with key change, observable effect, and average tool
call count.}
\label{tab:prompt-history}
\end{table}

\textbf{Notable observation.} Tool call count \emph{increased} from v1 to v4 because
agents were doing more useful work (checking tabs, validating state, recovering from
popups). The v5 budget cap reduced tool calls without reducing recovery because early
exit logic cut wasted calls on already-succeeded runs.

\section{Case studies focused on hard failures}

\subsection{Case A: nested iframe chain}

\textbf{Issue.} Primary host context showed a player container but the \texttt{src}
iframe pointed to a second-level page, which itself embedded a third-level player via
another iframe. The host agent harvested nothing because the actual HLS traffic
originated from the third-level frame.

\textbf{Resolution.} The host agent returned the second-level embed URL. The embedded
agent inspected the frame tree, detected the third-level origin, navigated directly
to it, clicked play, and successfully harvested the manifest.

\textbf{Lesson.} Nested iframe delegation is correctly handled by the hosting-to-embedded
fallback mechanism. Without this fallback, the run would have declared failure at the
host stage.

\subsection{Case B: popup and click trap}

\textbf{Issue.} The server button on the host page opened a full-screen ad popup each
time it was clicked. The popup occupied a new browser context. The original page
remained accessible in the background but the agent's focus was on the new context.

\textbf{Resolution.} The popup handler in v3 detected the new context via
\texttt{browser.pages()}, closed it, returned focus to the original page, and retried
the server click using \texttt{js\_click} (which bypasses the default click handler
that triggered the popup). The retry succeeded.

\subsection{Case C: tokenized stream generation}

\textbf{Issue.} No direct stream in initial page source. A decoy resource named
\texttt{mono.css} appeared in network traffic, which superficially resembled a
common pattern on a different provider. The actual signed m3u8 was generated
3 seconds after the play button was clicked.

\textbf{Resolution.} Post-activation harvest layers captured the runtime manifest URL.
The performance API scan retroactively found the m3u8 entry even though it had been
requested before \texttt{harvest()} was called.

\begin{figure}[H]
\centering
\begin{tikzpicture}[node distance=0.4cm and 0.8cm,
  case/.style={draw=accent!70, fill=blue!6, rounded corners=4pt,
               minimum width=4.5cm, minimum height=2.6cm,
               font=\small, align=center}
]

\node[case] (caseA) at (0,0) {
  \textbf{Case A: Nested iframe}\\[4pt]
  {\footnotesize Issue: 3-level frame chain}\\
  {\footnotesize Fix: embed fallback}\\
  {\footnotesize Result: manifest found}
};
\node[case] (caseB) at (5.2,0) {
  \textbf{Case B: Popup trap}\\[4pt]
  {\footnotesize Issue: ad popup on click}\\
  {\footnotesize Fix: close + js\_click}\\
  {\footnotesize Result: server activated}
};
\node[case] (caseC) at (10.4,0) {
  \textbf{Case C: Tokenized m3u8}\\[4pt]
  {\footnotesize Issue: no stream in HTML}\\
  {\footnotesize Fix: post-play harvest}\\
  {\footnotesize Result: manifest captured}
};

\node[flowbox, minimum width=13cm, fill=green!10, draw=green!60,
      below=1cm of caseB] (lesson)
  {All three cases resolved by the fallback + popup handler + deferred harvest chain};

\draw[myarrow] (caseA.south) -- ++(0,-0.3) -| (lesson.west);
\draw[myarrow] (caseB.south) -- (lesson.north);
\draw[myarrow] (caseC.south) -- ++(0,-0.3) -| (lesson.east);

\end{tikzpicture}
\caption{Hard-case recovery map: root cause, fix applied, and outcome for each case.}
\label{fig:case-map}
\end{figure}

\section{Cost and cache evaluation}

\subsection{Cache behavior}

Representative heavy-run pattern used for reporting (observed in late n8n-phase
campaigns and in the new Python implementation):

\begin{itemize}
  \item total context volume around \textasciitilde{}800k tokens;
  \item around \textasciitilde{}690k tokens served from cache;
  \item effective cache reuse around 80\%+ depending on accounting perspective and run
    phase.
\end{itemize}

\begin{table}[H]
\centering
\begin{tabularx}{\textwidth}{>{\raggedright\arraybackslash\bfseries}p{4cm}
                              >{\raggedright\arraybackslash}X
                              >{\raggedright\arraybackslash}p{2.5cm}}
\toprule
\rowcolor{accent!8}
Metric & Interpretation & Example value \\
\midrule
Total token context & aggregate context processed in heavy run scenarios & \textasciitilde{}800k \\
Cached token volume & portion served from cache layers & \textasciitilde{}690k \\
Cache effectiveness & practical reuse level across phases & \textasciitilde{}80\%+ \\
\bottomrule
\end{tabularx}
\caption{Representative cache and token pattern from project runs.}
\label{tab:cache-token}
\end{table}

\begin{table}[H]
\centering
\begin{tabularx}{\textwidth}{>{\raggedright\arraybackslash\bfseries}p{3.8cm}
                              >{\raggedright\arraybackslash}X
                              >{\raggedright\arraybackslash}X}
\toprule
\rowcolor{accent!8}
Scaling factor & Effect on runtime/cost & Mitigation \\
\midrule
More hosting pages & more specialist loops and tool calls &
  candidate ranking and pruning \\
More servers per page & more interaction + harvest cycles &
  bounded retries with confidence thresholds \\
More embedded fallbacks & deeper frame analysis cost &
  fallback gating and explicit stop criteria \\
\bottomrule
\end{tabularx}
\caption{Observed cost-scaling drivers in extraction campaigns.}
\label{tab:scaling}
\end{table}

\section{Overall evaluation summary}

Across all dimensions measured, the transition from intent-based to ReAct-style prompts
delivered the largest single improvement. Routing accuracy settled at around 95\%,
navigation recall went from 40\% to 89\%, and tokenized stream recovery went from roughly
25\% to 85\% by v3. Table~\ref{tab:prompt-history} in the prompt iteration section shows
the progression version by version.

\section{Discussion and observed strengths}

Looking at the evaluation results as a whole, a few things stand out. The explicit hosting-to-embedded fallback was probably the single most important architectural decision in terms of coverage: without it, anything involving a nested iframe chain just fails. Screenshot-grounded validation proved its worth too — before adding get\_media\_state and mandatory pre-harvest screenshots, it was impossible to tell from logs alone whether a server switch had actually done anything. The tool architecture also held up better than expected on unknown page structures; the multi-strategy interact fallback (selector, text, coordinates, js\_click) handled most selector drift without any prompt changes.

On the limitations side, the honest assessment is that anti-bot behavior is a moving target. Adblock filter lists needed refreshing multiple times over the evaluation period. Some Cloudflare-gated sessions still produce partial evidence that requires a human to judge. Tokenized stream timing is the most fragile piece — the 2-second post-play wait works for most providers but there are edge cases where the token exchange takes longer. And anything beyond a 3-level iframe chain hits the embedded agent's tool budget before it finds the actual player.

The prompt evolution from intent-based to ReAct-style was the most impactful change in terms of raw metrics, but it took several iterations to get the stopping criteria right. The budget cap in v3 was not obvious — it only became necessary after seeing agents burn through 20+ tool calls on paginated sites that had no additional content to find.

Chapter~7 draws on these evaluation results to summarize the project's contributions,
reflect on the methodology, and outline the short-term improvements and research
directions that would extend the platform toward production-grade, fully cloud-native
operation.
